\chapter{A Concrete Naming Certificate for $\CantorRealCompletionKernelUp$}
\label{ch:concrete-instances-cantor-real-completion-kernel-namecert}
\origin{human}

Following Cantor and Bishop--Bridges, the $\CantorRealCompletionKernelUp$ packet records the finite nested-dyadic completion route by which a decreasing interval chain becomes a regular rational name and then reaches the terminal $\RealUp$ seal. It sits after the dyadic endpoint ledger of \autoref{ch:concrete-instances-dyadicratcore-namecert}, the finite observation surface of \autoref{ch:concrete-instances-streamname-namecert}, the regular rational readback of \autoref{ch:concrete-instances-regseqrat-namecert}, the terminal real-name seal of \autoref{ch:concrete-instances-real-namecert}, and the nested dyadic interval carrier of \autoref{ch:concrete-instances-nesteddyadicinterval-namecert}. The packet names the completion kernel as displayed BEDC rows; it does not quotient streams, select a host limit, invoke countable choice, or assert ambient completeness for $\RealUp$.

\begin{definition}[Cantor real completion kernel carrier]
\label{def:cantor-real-completion-kernel-carrier}
A $\CantorRealCompletionKernelUp$ carrier is a finite $\BHist$ packet
$$
K=(I,D,W,Q,E,H,C,P,N)
$$
with the following displayed rows:
\begin{enumerate}
\item $I$ is the $\NestedDyadicIntervalUp$ chain row, listing the retained dyadic intervals and the successor containment ledger.
\item $D$ is the $\DyadicRatCoreUp$ endpoint and diameter row, recording lower endpoints, upper endpoints, midpoint reads, and dyadic diameter bounds for the retained windows.
\item $W$ is the $\StreamNameUp$ finite observation row through which a consumer requests inspected stages of the chain.
\item $Q$ is the $\RegSeqRatUp$ readback row extracted from the same windows and diameter bounds.
\item $E$ is the terminal $\RealUp$ seal row, available only after $I,D,W,Q$ have been displayed.
\item $H$ records componentwise $\hsame$ transport for interval rows, endpoint ledgers, observation windows, regular readbacks, terminal seals, replay, provenance, and local naming.
\item $C$ records the displayed $\Cont$ replay of the finite route from nested intervals to regular rational readback.
\item $P$ records
$$
\Pkg(K)=\Pkg(\CantorRealCompletionKernelUp,\NestedDyadicIntervalUp,\DyadicRatCoreUp,\StreamNameUp,\RegSeqRatUp,\RealUp,\BHist,\hsame,\Cont,\NameCert).
$$
\item $N$ is the local $\NameCert_{\CantorRealCompletionKernelUp}$ row.
\end{enumerate}
The classifier compares accepted packets only through $I,D,W,Q,E,H,C,P,N$. It has no coordinate for quotient real equality, a selected completed point, host interval intersection, countable choice, trichotomy, or ambient $\RealUp$ completeness.
\end{definition}

\begin{theorem}[Cantor kernel nested-containment obligations]
\label{thm:cantor-real-completion-kernel-nested-containment}
For every accepted $\CantorRealCompletionKernelUp$ carrier
$$
K=(I,D,W,Q,E,H,C,P,N),
$$
the local $\NameCert_{\CantorRealCompletionKernelUp}$ surface exposes successor containment only through the $\NestedDyadicIntervalUp$ row $I$ and the endpoint comparisons recorded in the $\DyadicRatCoreUp$ row $D$. Every displayed adjacent pair in the retained chain must be read as a finite containment certificate before the corresponding $\StreamNameUp$ window, $\RegSeqRatUp$ readback, or $\RealUp$ seal can be consumed.
\end{theorem}
\begin{proof}
Project the carrier of \autoref{def:cantor-real-completion-kernel-carrier}. The rows available before any regular readback are the nested interval row $I$, the dyadic endpoint row $D$, and the structural transport rows.

Read the successor containment ledger in $I$. Each adjacent retained pair is paired with lower-endpoint and upper-endpoint comparisons in $D$, so containment is witnessed by finite dyadic data rather than by a completed intersection point.

Only after those rows are displayed may the consumer inspect $W,Q,E$. Since the carrier has no quotient, selected limit, host intersection, or ambient-completeness coordinate, the containment obligation is exactly the finite route stated above.
\end{proof}

\begin{theorem}[Cantor kernel shrinking-diameter obligations]
\label{thm:cantor-real-completion-kernel-shrinking-diameter}
For every accepted $\CantorRealCompletionKernelUp$ carrier
$$
K=(I,D,W,Q,E,H,C,P,N),
$$
each precision request read by the $\StreamNameUp$ row $W$ must be matched by a displayed dyadic diameter bound in $D$ for a retained interval of $I$. The $\RegSeqRatUp$ row $Q$ may read a Cauchy modulus only through those bounds, and no consumer can replace the displayed shrinking ledger by an ambient compactness or completeness principle.
\end{theorem}
\begin{proof}
Begin with the finite observation row $W$. It names only the inspected stages of the nested dyadic chain, so a precision request must point back to one retained interval row of $I$.

The dyadic row $D$ supplies the corresponding diameter bound. That bound is a displayed dyadic comparison attached to the same retained interval, not a host statement about all unseen stages.

The readback row $Q$ consumes the pair $(W,D)$ as regular rational completion data. The packet therefore exposes shrinking diameter through finite windows and dyadic bounds, while the excluded coordinates prevent compactness, selected limits, quotient equality, or ambient completeness from entering the route.
\end{proof}

\begin{theorem}[Cantor kernel RegSeqRat-to-Real sealing]
\label{thm:cantor-real-completion-kernel-regseqrat-real-sealing}
Every $\RealUp$-facing consumer of an accepted $\CantorRealCompletionKernelUp$ carrier must pass through the displayed route
$$
\NestedDyadicIntervalUp \longmapsto \DyadicRatCoreUp \longmapsto \StreamNameUp \longmapsto \RegSeqRatUp \longmapsto \RealUp
$$
using the rows $I,D,W,Q,E,H,C,P,N$. The terminal seal $E$ is not readable until the nested-containment and shrinking-diameter obligations have been displayed through $I,D,W,Q$.
\end{theorem}
\begin{proof}
Apply \autoref{thm:cantor-real-completion-kernel-nested-containment}. It forces the consumer to read the nested interval row $I$ and the dyadic endpoint comparisons in $D$ before treating the chain as a valid decreasing dyadic packet.

Apply \autoref{thm:cantor-real-completion-kernel-shrinking-diameter}. It forces every precision request in $W$ to be witnessed by a displayed dyadic diameter bound in $D$, and it routes the resulting Cauchy data through $Q$.

Finally read the terminal seal row $E$ through $H,C,P,N$. These structural rows transport, replay, source, and name only the displayed route. Since the packet contains no selected point, quotient real equality, host limit, countable-choice schedule, or ambient completeness coordinate, the $\RealUp$ seal is reached only by the stated $\RegSeqRatUp$ handoff.
\end{proof}

\falsifiablePrediction{An accepted $\CantorRealCompletionKernelUp$ consumer is refuted if it reaches a $\RealUp$ seal without displaying the nested dyadic interval row, dyadic endpoint and diameter ledger, finite StreamName windows, RegSeqRat readback, componentwise hsame transport, Cont replay, Pkg provenance, and local NameCert row.}

\closureat{CantorRealCompletionKernelUp}{seedStr}

\begin{closurestatus}{\CantorRealCompletionKernelUp}
  \constructivestory{A $\CantorRealCompletionKernelUp$ packet is generated from a NestedDyadicInterval chain row, a DyadicRatCore endpoint and diameter ledger, StreamName finite windows, RegSeqRat readback, a RealUp seal, componentwise hsame transport, Cont replay, Pkg provenance, and a local NameCert row.}
  \theoryclosure{\seedClosure}
  \scopeclosed{\autoref{def:cantor-real-completion-kernel-carrier}, \autoref{thm:cantor-real-completion-kernel-nested-containment}, \autoref{thm:cantor-real-completion-kernel-shrinking-diameter}, and \autoref{thm:cantor-real-completion-kernel-regseqrat-real-sealing} seed the finite Cantor completion route from nested dyadic intervals through dyadic endpoint and diameter rows, finite windows, regular rational readback, terminal Real sealing, transport, replay, provenance, and local naming.}
  \formalstatus{\unformalizedV}
  \bridgestatus{none}
  \notclaimed{This chapter does not close quotient real equality, selected host limits, countable choice, trichotomy, ambient $\RealUp$ completeness, Lean verification, axiom-clean audit, or bridge checking.}
  \upgradepath{Obligation closure requires checked carrier admission, nested containment, shrinking diameter, RegSeqRat readback, Real seal nonescape, transport, replay, provenance, and local naming for BEDC.Derived.CantorRealCompletionKernelUp.}
\end{closurestatus}
